\section{Accuracy Does Not Transfer to the Queue}
\label{sec:accuracy}

Section~\ref{sec:ordering} found little headroom above the score. We test that
directly by holding the score and its weights fixed and swapping the classifier
underneath it, replacing the LaBSE priority and sentiment posteriors with those
of the TF-IDF baseline (Table~\ref{tab:subst}).

\begin{table}[!tb]
\caption{Same score, same weights, posteriors swapped. $\rho=1.05$, 200 seeds,
waits in minutes for gold-High tickets.}
\label{tab:subst}
\centering
\small
\begin{tabular}{@{}lrr@{}}
\toprule
 & LaBSE & TF-IDF \\
\midrule
priority \macrof{} (test) & \textbf{0.8901} & 0.8706 \\
\midrule
mean wait      & 1.98 & \textbf{1.80} \\
median wait    & 1.11 & \textbf{1.10} \\
p95 wait       & \textbf{5.32} & 5.42 \\
worst-case wait & 362.1 & \textbf{97.9} \\
\midrule
gold-High missed & \textbf{165} & 193 \\
\quad of which buried below 0.1 & 68 & \textbf{8} \\
\bottomrule
\end{tabular}
\end{table}

Through the median and the 95th percentile the two are indistinguishable, at
tenths of a minute and with the sign changing between rows. They separate in the
tail, where the more accurate model's worst-served urgent ticket waits
\textbf{six hours against the weaker model's ninety-eight minutes}.

The claim is not that TF-IDF wins. It is that a two-point \macrof{} advantage
buys nothing through the body of the queue and costs something at the extreme,
which no classification table would report.

\textbf{The mechanism is a metric F\textsubscript{1} cannot see.} LaBSE misses
fewer urgent tickets than TF-IDF, 165 against 193, and buries far more of the
ones it misses: 68 fall below a score of 0.1 against TF-IDF's 8. Fewer mistakes,
$8.5\times$ more confident ones. A confidently-missed urgent ticket falls to the
bottom of the queue and stays there, which is the six-hour worst case above,
while a hedged miss lands mid-queue and is picked up soon. Macro-F\textsubscript{1}
counts both as one error.

We therefore propose \textbf{buried urgent tickets}, the count of gold-High
tickets scored below a low threshold, as a triage-specific error metric worth
reporting beside F\textsubscript{1}. It is what the queue actually feels.

\subsection{A hypothesis we refuted}

We expected calibration to explain the disparity and to fix it cheaply. The
diagnosis supported it: LaBSE is more accurate than TF-IDF, $3.7\times$ worse
calibrated (ECE 0.0658 against 0.0180), and worst calibrated on exactly the track
it is least accurate on (Tanglish, 0.1071). That is a double penalty, and
temperature scaling \cite{guo2017calibration} is the standard remedy.

Fitting the temperature on a clean held-out development set fixed the
calibration and made the queue \emph{worse}. Mean wait for urgent tickets rose
from 1.98 to 2.65 minutes and the Tanglish-minus-English gap widened from $+0.50$
to $+0.97$. The reason is that expected calibration error measures agreement
between top-1 confidence and accuracy, while a queue needs posteriors that
\emph{separate} tickets, and temperature scaling compresses exactly that
separation. "Calibrate before you score" was wrong here, and we report it because
the diagnosis is convincing enough that others will try it.
